\chapter{参考资料}\label{app:references}
本书的结构、文字和示例代码重新编写。以下书籍用于确定学习范围；书名与版次按所持资料标注，缺少版权页信息的条目不补猜测的出版年。涉及库的接口行为，以本书锁定的依赖版本和官方文档为准；网页资料核对日期为 2026 年 11 月 5 日。

\section*{书籍}
\begin{enumerate}
\item 崔庆才：《Python 3 网络爬虫开发实战（第二版）》；版次见书名。
\item 李玺：《爬虫逆向进阶实战》。
\item 李岳阳、卓斌：《反爬虫 AST 原理与还原混淆实战》。
\item Dimitrios Kouzis-Loukas：《精通 Python 爬虫框架 Scrapy》。
\item Ryan Mitchell：《Python 网络数据采集》。
\end{enumerate}

\section*{官方文档与标准}
\begin{enumerate}
\item Python 3.12 标准库：\href{https://docs.python.org/3.12/library/html.parser.html}{HTMLParser}、\href{https://docs.python.org/3.12/library/urllib.parse.html}{URL 解析}、\href{https://docs.python.org/3.12/library/sqlite3.html}{SQLite 与连接生命周期}、\href{https://docs.python.org/3.12/library/csv.html}{CSV}、\href{https://docs.python.org/3.12/library/asyncio.html}{asyncio}。
\item \href{https://requests.readthedocs.io/en/stable/user/quickstart/}{Requests 官方快速入门}；\href{https://www.crummy.com/software/BeautifulSoup/bs4/doc/}{Beautiful Soup 官方文档}。
\item \href{https://docs.scrapy.org/en/latest/intro/tutorial.html}{Scrapy 官方教程}与\href{https://docs.scrapy.org/en/latest/topics/item-pipeline.html}{Item Pipeline}；\href{https://playwright.dev/python/docs/intro}{Playwright Python 官方文档}。
\item IETF：\href{https://www.rfc-editor.org/rfc/rfc9309.html}{RFC 9309，Robots Exclusion Protocol}。
\item \href{https://developer.chrome.com/docs/devtools/network}{Chrome DevTools 网络面板}；\href{https://developer.android.com/studio/debug/network-profiler}{Android Studio Network Inspector}。
\item Babel：\href{https://babeljs.io/docs/babel-parser}{Parser}、\href{https://babeljs.io/docs/babel-traverse}{Traverse} 与 \href{https://babeljs.io/docs/babel-generator}{Generator}。
\item \href{https://lxml.de/xpathxslt.html}{lxml 的 XPath 文档}；\href{https://docs.aiohttp.org/en/stable/client_quickstart.html}{aiohttp 客户端文档}；\href{https://docs.scrapy.org/en/latest/topics/item-pipeline.html}{Scrapy Item Pipeline 文档}。
\item \href{https://urllib3.readthedocs.io/en/stable/reference/urllib3.util.html}{urllib3 Retry 文档}；\href{https://requests.readthedocs.io/en/latest/user/quickstart/}{Requests 编码、超时与重定向说明}。
\end{enumerate}
